\documentclass{report}
\usepackage{amsmath,amssymb}
\setlength{\parindent}{0mm}
\setlength{\parskip}{1em}
\begin{document}
\begin{center}
\rule{6in}{1pt} \
{\large Andrew, R Reisner \\
{\bf Progress on Improving the Parallel Scalability of BoxMG}}

201 N Goodwin Ave \\ Urbana \\ IL \\ 61801
\\
{\tt areisne2@illinois.edu}\\
Luke Olson\\
David Moulton\end{center}

Black Box Multigrid (BoxMG) is a variational multigrid method for discretizations
of PDEs on logically structured grids.
BoxMG uses operator-induced interpolation, making it robust for problems
with discontinuous coefficients and general boundary conditions on grids
of any dimension (i.e., non-power-of-two grids).
Like most structured multigrid methods, BoxMG uses the structure of the problem to
yield notable performance advantages such as direct memory access and bounded
complexity in the coarse grid operator. Yet, while current parallel BoxMG
implementations scale well out to thousands of cores,
the coarse grid solve and line-relaxation limit further scalability.
In the plasma applications of interest, grid stretching creates an
anisotropic diffusion operator, and hence, with standard coarsening
line-relaxation is required.
The current coarse-grid solve strategy first collects the coarse problem to a single
core. The coarse problem is then solved using serial BoxMG.
To extend the scalability of the coarse-grid solve, we look at redistributing
the problem on coarser levels to provide redundancy while minimizing communication costs.
In this talk, we give an overview of the current state of BoxMG including a new control
layer. We will show our progress in addressing the coarse-grid scalability limitations,
specifically in the context of electric field calculations on curvilinear meshes.


\end{document}
